\section{Analysis Challenges}\label{sec:Challenges}

Designing a static analysis to detect missing \texttt{RB\_GC\_GUARD} directives presents significant technical challenges. This section characterizes the fundamental difficulties that any such analysis must address, which inform both our choice of analysis framework and the design of our detection algorithm.

\subsection{Required Analysis Capabilities}

To detect omitted guards, an analyzer must answer four fundamental questions about the program. First, the analyzer must identify all program points where a raw C pointer is derived from a Ruby \texttt{VALUE}. These derivation sites include macro invocations such as \texttt{RSTRING\_PTR}, and \texttt{TypedData\_Get\_Struct}, as well as function calls that return pointers into object-owned memory. Each derivation establishes a dependency relationship: the validity of the subordinate pointer depends on the liveness of the owning \texttt{VALUE}.

Second, the analyzer must identify call sites that may invoke the garbage collector, either directly or transitively. Any such call site represents a potential hazard point where, if the owning \texttt{VALUE} has become invisible to conservative scanning, the subordinate pointer may become dangling.

Third, the analyzer must identify uses of subordinate pointers that occur after a potential GC trigger in the control flow. A ``use'' includes not only direct dereferences but also passing the pointer as an argument to other functions, which may dereference it internally.

Fourth, the analyzer must determine whether the owning \texttt{VALUE} remains visible to the GC between the trigger and the use of the subordinate pointer. Protection can come from two sources: an explicit \texttt{RB\_GC\_GUARD} directive placed after the trigger, or a subsequent \texttt{VALUE}-level access that forces the compiler to keep it live. When neither protection exists, a missing guard should be flagged.

Combining these requirements, a missing guard is indicated when a subordinate pointer $p$ is obtained from a \texttt{VALUE} $v$, a GC-triggering call occurs after the derivation, the pointer $p$ is used after the trigger, and the variable $v$ has no \texttt{VALUE}-level access after the trigger and no explicit guard.

\subsection{Scale of the Codebase}

CRuby comprises over 300,000 lines of C code with more than 23,000 \texttt{VALUE} variables and over 42,000 functions. Naive enumeration of all possible derivation-trigger-use triples would be computationally impractical: a cubic algorithm over these quantities could require examining on the order of $10^{13}$ combinations. Any practical analyzer must employ efficient indexing and pruning strategies to focus on relevant program fragments while completing in a reasonable time frame.

\subsection{Pervasiveness of GC Triggers}

The Ruby runtime is deeply interconnected, and most non-trivial operations can transitively reach the GC entry point. In CRuby 3.4.5, our analysis found that approximately 25\% of all functions (10,608 out of 42,612) can reach \texttt{gc\_enter}. This high proportion arises because object allocation (\texttt{rb\_str\_new}, \texttt{rb\_ary\_new}, etc.) may trigger GC when the heap is full, method dispatch (\texttt{rb\_funcall}) may allocate during argument preparation or result conversion, exception handling paths often allocate exception objects, and string operations (\texttt{rb\_str\_cat}, \texttt{rb\_str\_concat}) may resize buffers, triggering allocation.

This means that distinguishing ``safe'' calls from ``potentially triggering'' calls requires whole-program analysis rather than local inspection. A simple heuristic (e.g., ``only \texttt{rb\_gc} triggers GC'') would miss the vast majority of real hazards. Identifying all functions that may transitively invoke GC requires computing the transitive closure of the call graph rooted at GC entry points---a task that demands efficient recursive query evaluation over a large call graph.

\subsection{Diversity of Derivation Patterns}

Subordinate pointers are obtained through many different APIs and macros, each with different signatures and expansion patterns. Some are simple macro expansions (e.g., \texttt{RSTRING\_PTR}), while others involve function calls with complex control flow. Table~\ref{tab:derivation_patterns} summarizes the major categories.

\begin{table}[ht]
\centering
\caption{Categories of derivation patterns in CRuby}
\label{tab:derivation_patterns}
\begin{tabular}{lll}
\toprule
\textbf{Category} & \textbf{Examples} & \textbf{Return Type} \\
\midrule
String access & \texttt{RSTRING\_PTR}, \texttt{StringValueCStr} & \texttt{char *} \\
Array access & \texttt{RARRAY\_PTR}, \texttt{RARRAY\_CONST\_PTR} & \texttt{VALUE *} \\
Typed data & \texttt{TypedData\_Get\_Struct}, \texttt{DATA\_PTR} & \texttt{void *} \\
Binary data & \texttt{rb\_io\_buffer\_get\_bytes} & \texttt{void *} \\
\bottomrule
\end{tabular}
\end{table}

Additionally, subordinate pointers may be passed through multiple variables or returned from helper functions, requiring tracking of where pointer values originate. Full interprocedural tracking of pointer origins is expensive and may introduce additional imprecision.

\subsection{Detecting Existing Guards}

The \texttt{RB\_GC\_GUARD} mechanism is implemented as a macro that expands differently across compilers and configurations. Listing~\ref{lst:guard_expansions} shows how the same source-level guard can appear differently in preprocessed code.

\begin{figure}[tb]
\begin{lstlisting}[caption={Different expansions of \texttt{RB\_GC\_GUARD} across configurations.},label={lst:guard_expansions}]
// Source code:
RB_GC_GUARD(str);

// GNU expansion (typical):
(*__extension__ ({
    volatile VALUE *rb_gc_guarded_ptr = &(str);
    __asm__("" : : "m"(rb_gc_guarded_ptr));
    rb_gc_guarded_ptr;
}));

// MSVC expansion (alternative):
(*rb_gc_guarded_ptr(&(str)));
\end{lstlisting}
\end{figure}

Reliably detecting whether a guard is present requires either access to macro invocation metadata or recognition of characteristic expansion patterns in the preprocessed code. The analyzer must handle both cases to be robust across build configurations.

\subsection{Precision-Recall Tradeoff}

The analysis must navigate a precision-recall tradeoff. Over-approximating derivations or triggers leads to excessive false positives, while under-approximating them risks missing real bugs. Given that CRuby has been maintained for 30 years with extensive testing, the analyzer must be precise enough to surface actionable candidates rather than overwhelming maintainers with noise.

Table~\ref{tab:precision_tradeoffs} summarizes the key tradeoffs in our design.

\begin{table}[ht]
\centering
\caption{Precision-recall tradeoffs in analysis design}
\label{tab:precision_tradeoffs}
\begin{tabular}{p{2.5cm}p{2.5cm}p{2.5cm}}
\toprule
\textbf{Component} & \textbf{Over-approx.\ effect} & \textbf{Under-approx.\ effect} \\
\midrule
Derivations & False positives from non-hazardous pointers & Missed bugs from unmodeled APIs \\
GC triggers & False positives from safe calls & Missed bugs from indirect triggers \\
Pointer uses & False positives from dead code paths & Missed bugs from unrecognized uses \\
\bottomrule
\end{tabular}
\end{table}
